\clearpage
\section*{2026-09-16: Five more priority sites clear the land review}
\textit{Agent-drafted research entry.}

Recorded documents establish physical building parcels for eight of the nine
remaining post-policy parents with an exact-99 constituent. Five also clear the
combined land-and-filing screen: North 8th Street, Shepherd/Highland, East 108th
Street, East 178th Street, and Beach 29/30. The unresolved count falls from 158 to 153 of
877 weighting-eligible parents: 119 historical and 34 post-policy. The audit
now has 241 candidate area corrections among 724 passing parents. All
production observations, memberships, dates, units, areas, and weights remain
unchanged during this source review.

The supported North 8th building parcel is approximately 7,238 square feet,
compared with 12,450 in production. Its May 2026 development agreement assigns
lot 133 to the 99-unit ground lease and retains the church on the neighboring
parcels. Recorded legal courses supply the area. At East 108th, the legal rectangles for the two-building developer
parcel total 44,200 square feet, compared with 77,600 in production. Retained
neighboring lot 342 provides nonexclusive parking/access rights but remains
separate ground. The two mapped building lots later split again to create
lot 339, which lies inside that developer parcel.

Shepherd/Highland retains its full 20,653-square-foot earlier parcel. The
December 2025 diagram and June 2026 mortgage schedules assign both lots 48
and 148 to Highland and lot 49 to Shepherd. A September 2026 filing proposes
21 units on lot 148, after the July 8 sample cutoff; it does not add units to
the observed 198-unit pair. East 178th's live replacement occupies former
lots 88--90, whose archival areas total 7,287 square feet. The earlier proposed
lot 87 is outside its recorded developer parcel. Filing dates remain unchanged.

Three other priority cases have supported ground but retain filing questions.
Charles Place's developer land is merged lot 71, formerly lots 7 and 71,
measuring 12,590 square feet. Broadway's building parcel remains 9,262 square
feet, with adjacent lot 28 providing an easement rather than added ground.
Their older 46-unit and hotel/residential filings respectively need reconciliation
with the current 99-unit proposals. Godwin/Kimberly's July agreement explicitly
assigns 99 units to lot 79 and 65 to lot 88 under the same owner. Together with
the common March 31, 2025 filing date and architect, this supports a proposed
164-unit economic parent. The merger is a supervisory recommendation pending
implementation and downstream verification. Its physical building parcels
measure 12,541 and 5,900 square feet. Old lot 78's 17,670-square-foot attribute
is inconsistent with its roughly 28,483-square-foot mapped extent and must
not be allocated proportionally.

The 165th Street condominium maps describe the earlier commercial site; the
2025 six-building residential plan remains necessary. Beach 29/30 is now
resolved by recorded agreements: developer lots 44/48 and 50 total approximately
47,365 square feet. Fifteen narrow parcels retain their own 30-unit allocation
and independent lot coverage. The fifteen small-house filings match those
retained parcels, so they do not imply shared apartment ground. The code
preserves those filings and clears only their documented land conflicts.
Broader zoning lots remain separate from physical development land.

The Sol inventories cover all 158 starting cases: 62 in Brooklyn, 48 in the
Bronx, and 48 elsewhere. Most still require primary-source adjudication.
Additional missed-link candidates are GO Broome's Suffolk/Norfolk buildings
and the Sutphin/95th Avenue pair. Rockaway Village's three same-day filings
remain one parent; separate financing phases alone do not warrant a split.

\textit{Reproduction.} Working changes based on \texttt{a665b53}; run
\texttt{make dof-site-review}. The audit's eleven-row
\path{parent_site_document_review.csv} validates archival parcels against
MapPLUTO 23v3\_1 and retains independent scope flags. The supervised report
records document IDs, pages, and area calculations. Twenty-one original ACRIS
PDFs, including sixteen recovered here, have source URLs and SHA-256 hashes.
